\documentclass[]{../../../../tex/classes/styledArticle}
\usepackage{../../../../tex/packages/hebrewSupport}
\usepackage{../../../../tex/packages/mathShortcuts}
\usepackage{../../../../tex/packages/theoremsSupport}

\author{שחר פרץ}
\title{חדו''א 2א $\sim$ \textit{הרצאה 11}}
\begin{document}
	\maketitle
	\textbf{מרצה: }יעל קרשון
	\begin{Exercise}[תרגיל פתיחה]
		\begin{enumerate}
			\item נתונות $u_n \co [0, R] \to \R$ רציפות. נניח שהטור $\sumnoinf u_n(x)$ מתכנסת במ''ש ב־$[0, R)$. אז $\sumnoinf u_n(x)$ מתכנס במ''ש ב־$[0, R]$. 
			\item ידוע ש־$B_n(x) := \sumnko \sin(kx) = \frac{\sin\cl{\frac{n + 1}{2}x}\sin\cl{\frac{n}{2}x}}{\sin\cl{\frac{x}{2}}}$ (ברשימות). הראנו ש־$\{B_n\}$ חסומה במידה אחידה ב־$[\eg, 2\pi - \eg]$. הוכיחו שהמשפחה $\{B_n\}$ אינה חסומה במ''ש ב־$(0, 2\pi)$. 
			
			\textbf{רמז: }מצאו $x_n \in (0, 2\pi)$ כך ש־$x_n \to 0$ ולכל $n$ מתקיים $\sin\cl{\frac{n}{2}x_n} = 1$. הראו ש־$B_n(x_n) \toinf \infty$. 
		\end{enumerate}
	\end{Exercise}
	
	\subsection*{פונקציה רציפה שאינה גזירה באף נקודה}
	ממליץ מאוד לראות של הגרפים ברשימות של שירי: (נסמן ב־$\{x\} := x - \floor{x}$ את החלק השברי של $x$)
	\begin{align*}
		u_0(x) := \begin{cases}
			\{x\} & n \le x \le n + \frac{1}{2} \\
			1 - \{x\} & n + \frac{1}{2} \le x \le n + 1
		\end{cases} && u_k(x) := \frac{1}{4^{k}}u_0(4^{k}x)
	\end{align*}
	ניעזר בפונקציות העזר האלו, כדי להתבונן בטור $f(x) := \sumkoinf u_k(x)$ שמתכנס במ''ש. זה מבוחן $M$ של ווייראשטראס: עבור $M_k = \frac{1}{24^{k}}$ נקבל $\sof{u_k(x)} \le M_k$, ו־$\sumkoinf M_k$ מתכנס. מכאן נובע ש־$f$ מוגדרת היטב ורציפה (התכנסות במ''ש של רציפות היא רציפה). יתר על כן, $f$ איננה גזירה באף נקודה – נוכיח זאת. 
	
	\theo{לכל $x_0 \in \R$ קיימת סדרה $x_n$ כך ש־
	\[ \frac{f(x_n) - f(x_0)}{x_n - x_0} \in \begin{cases}
		\Neven & n \in \Neven \\
		\Nodd & n \in \Nodd
	\end{cases} \]}
	ההוכחה של המשפט ברשימות של שירי. בפרט, נובע מכך שלא קיים גבול של $\lim_{\eta \to x_0}\frac{f(x_n) - f(\eta)}{x_n - \eta}$ מהיינה. 
	
	עד עתה, סיימנו את \shiri{4.3}. נמשיך ל־\shiri{4.4} – הסעיף האחרון לפני שנעבור לדבר על טורי חזקות. 
	\subsection*{קירוב במ''ש ע''י פולינומים}
	\textbf{תזכורת: }
	\begin{itemize}
		\item $\sumnoinf \frac{x^{n}}{n = 0} = \exp(x)$ במ''ש על קטעים סגורים ב־$\R$. 
		\item $\sumnoinf x^{n} = \frac{1}{1- x}$ במ''ש על קטעי־קטעים סגורים של $(-1, 1)$. 
	\end{itemize}
	הטורים לעיל הם דוגמאות ל\textit{טורי חזקות}, שזה הנושא הבא. בפרט, בכל קטע סגור $[a, b]$, קיימת סדרה של פולינומים $P_n(x) = \sumnko \frac{x^{k}}{k!}$ כך ש־$P_n \to \exp(x)$ במ''ש. 
	
	''אוי, שכחתי לקטום``
	
	זה מקרה מיוחד, שבו כל $P_n$ נוצר מקודמו ע''י חיבור של פולינום ממעלה גבוהה יותר. אם מתירים להשתמש בפולינומים כלליים, מגיעים למשפט וויראשטראס. 
	\begin{Theorem}[וויראטשראס]
		תהי $f \co C([a, b])$ כלשהי. אז קיימת סדרה של פולינומים $P_k(x)$ כך ש־$P_k \to f$ במ''ש\footnote{מומלץ לקרוא על ה־function topology של $C([a, b])$, ואז משפט וויראשטראס למעשה אומר צפיפות של הפולינומים בתוך מרחב הפונקציות הרציפות. }. 
	\end{Theorem}
	\rmark{בד''כ \textit{אין} טור חזקות ש־$P_k$ סכומיו החלקיים (יש פרשנות פורמלית לטענה הזו). ''אבל אם אתם פוגשים פונקציה בתחנת האוטובוס היא כנראה סימפטית מספיק``. פונקציות שניתנות ע''י טור חזקות נקראות \textit{פונקציות אנליטיות}, ויש להן תכונות מאוד נחמדות. }
	אומנם, כל פולינומים אפשר לבטא כסכום של פולינומים. אך המקדמים תלויים ב־$P_k$, כלומר $P_k(x) = \sum_{j = 0}^{d_k}c_{j, k}x^{j}$. 
	
	אם $[a, b] = [0, 1]$, אז $B_n(x) = f(x)$ במ''ש כאשר: 
	\[ B_n(x) = \sumnko \binom{n}{k}f\cl{\frac{n}{k}}x^{k}(1 - x)^{n - k} \]
	שימו לב שכאן המקדמים משתנים. בכל מקרה, הוכחה אצל שירי. 
	
	\section*{טורי פונקציות}
	ברשימות של שירי, \shiri{5}. בשיעור הזה נשתדל להספיק את \shiri{5.1}-\shiri{5.4}. 
	\defi{\textit{טור חזקות סביב $x = \eta$} הוא טור פורמלי (כלומר אכפת לנו מהמקדמים בלבד) מהצורה $\sumkoinf a_k(x - \eta)^{k}$}
	''אטא זה פשוט קשקוש אחר`` – סטודנטית אחרי שיעל החליפה $\xi$ ב־$\eta$. 
	
	לעת עתה, לצורך הנוחות, נניח בה''כ ש־$\eta = 0$, אז ברור שהטור מתכנס ב־$x = 0$. אבל מה קורה לגבי שאר ה־$x$־ים? 
	\lem{נניח $\sum a_k x^{k}$ מתכנס ב־$x_0 \neq 0$. יהי $R := \sof{x_0}$. יהי $0 < R' < R$, אז הטור מתכנס במ''ש בקטע $[-R', R']$. }
	\rmark{שימו לב, הטור לאו דווקא מתכנס בקטע הפתוח $(-R, R)$. }\begin{proof}
		לכל $x \in [-R', R']$ נרשום $\sof{a_kx^{k}} = \sof{a_kx^{k}}\cdot \sof{\frac{x}{x_0}}^{k}$. הטור $\sum a_kx^{k}$ מתכנס, ולכן סדרת איבריו שואפת לאפס, ולכן חסומה. יהי $B := \sup_{k \in \Z}\sof{a_kx_0^{k}}$. תהי $q:=\frac{R'}{R}$. אזי $0 < q < 1$, ולכן $\sumkoinf Bq^{k}$ מתכנס. עתה, לכל $x \in [-R', R']$ מתקיים: 
		\[ \sof{a_kx^{k}} = \underbrace{\sof{a_kx_0^{k}}}_{\le B}\underbrace{\sof{\frac{x}{x_0}}^{k}}_{\mathclap{=\frac{\sof x ^{k}}{R^{k}} \le q^{k}}} \le Bq^{k} \]
		ולכן מבוחן $M$ של וויראשטראס $\sum a_kx^{k}$ מתכנס בהחלט במ''ש ב־$[-R', R']$. 
	\end{proof}
	
	בשלב הזה יהיו לנו כל מיני מסקנות נחמדות. 
	\cola{יהי נתון טור חזקות $\sumkoinf a_kx^{k}$. אז: 
	\begin{enumerate}[(A)]
		\item הטור מתכנס ב־$(-R, R)$. נוכל להגדיר $f \co (-R, R) \to \R$ ע''י $f(x) = \sumkoinf a_kx^{k}$. מעתה ואילך, נגדיר $f(x) := \sumkoinf a_kx^{k}$. 
		\item $f$ רציפה. 
		\item גם הטור $\sumkoinf a_k\frac{x^{k + 1}}{k + 1}$ מתכנס ב־$(-R, R)$, ולכל $x \in (-R, R)$ מתקיים: 
		\[ \int_0^{x}f(t)\dt = \sumkoinf \frac{x^{k}}{k + 1} \]
	\end{enumerate}}
	\begin{proof}
		\begin{enumerate}[(A)]
			\item יהי $\hat x \in (-R, R)$. יהי $\sof{\hat x} < R' < R$. אז $-R < -R' < \hat x < R ' < R$. מהלמה שהוכחנו, הטור מתכנס במ''ש ב־$[-R', R']$, ובפרט מתכנס נקודתית ב־$\hat x$. 
			\item $f|_{[-R', R']}$ גבול במ''ש של רציפות, לכן $f$ רציפה ב־$[-R', R']$. בפרט, $f$ רציפה ב־$\hat x$. 
			\item ב־$[-R', R']$ הטור מתכנס במ''ש, ולכן מהמשפט הנוגע לאינטגרציה איבר־איבר, נקבל ש־
			\[ \int^{x}_0 \dequad \ f(t)\dt = \sumkoinf a_k \cdot \frac{x^{k}}{k + 1} \]
			במ''ש ב־$x \in [-R', R']$. בפרט ב־$x = \hat x$ (כי $\int x_k = \frac{x^{k + 1}}{k + 1}$)
		\end{enumerate}
	\end{proof}
	
	את שני המשפטים הבאים, משפט אבל ומשפט קושי־הדמרד, הוכחנו בחדו''א 1א. נבצע הוכחות דומות היום. 
	
	\begin{Theorem}[אבל]
		יהי טור חזקות $\sumkoinf a_kx^{k}$. אז קיים $0 \le R \le \infty$, כך ש־: 
		\begin{itemize}
			\item אם $\sof{x} < R$ אז הטור מתכנס ב־$x$. 
			\item אם $\sof x > R$ אז הטור מתבדר ב־$x$. 
			\item אם $\sof x = R$ אז הולכים לשבור את הראש מה קורה ב־boundary. 
		\end{itemize}
		ה־$R$ הזה נקרא \textit{רדיוס התכנסות}\footnote{במרוכבות, זה אותו הדבר אבל אשכרה רדיוס, כי מדובר על דיסק. גם במקרה המרוכב, צריך לבדוק ידנית מה קורה ב־boundary (כלומר ב־$\partial\{x \le R \mid x \in \R\} \cong S^{1}$)}. 
	\end{Theorem}
	זה משפט שאנו מכירים מחדו''א 1א. נציג הוכחה דומה. 
	\begin{proof}
		נגדיר: 
		\[ R = \sup\ccb{\sof{x_0} \co  \text{מתכנס}\ \sum a_n x_0^{n}} \]
		נניח $\sof x < R$. מהגדרת הסופרמום קיים $x_0$ כך ש־$\sof x < \sof{x_0} < R$. מהלמה נובע ש־$\sumninf a_nx^{n}$ מתכנס. 
		
		במקרה השני, נניח $\sof x > R$. מהגדרת הסופרמום, $\sum a_nx^{k}$ מתבדר. 
	\end{proof}
	\rmark{''למי שרואה את הסרטון שבוע לפני הבחינה – בבוחן ראינו שיש בעיות מהותיות``. יעל מעודדת אותנו לחזור על חומרים של חדו''א 1א. }
	\rmark{ה־$R$ בהוכחה לעיל מוגדר, משום שהקבוצה לא ריקה כי הטור מתכנס בהכרח ב־$x = 0$. }
	\begin{Theorem}[קושי־הדמרד]
		יהי $R$ רדיוס ההתכנסות של $\sumkoinf \sof{a_k}x^{k}$. אז: 
		\[ \frac{1}{R} = \limsup_{k \to \infty}\sqrt[k]{a_k} \]
		כאשר $R$ רדיוס ההתכנסות, במובן הרחב (אם ה־$\limsup$ הוא $0$, הטור מתכנס בכל $\R$, אם הוא מספר אז במובן הצר, ואם הוא $\infty$ אז הוא מתכנס בסינגילטון $\{x_0\}$)
	\end{Theorem}
	\begin{proof}
		נגדיר $\bar{R} := \cl{\limsup_{k \to \infty}\sof{a_k}^{\frac{1}{k}}}\op$ במובן הרחב. יהי $x \in \R$. נניח $0 < \sof x < \bar R$ (במובן הרחב). אז ישנו $q \in (0, 1)$ כך ש־$\sof x < q \bar R$. מהגדרת $\bar R$ מתקיים $\limsup \sof{a_k}^{\frac{1}{k}} < \frac{q}{\sof x}$. 
		
		תרגיל: קיים $k_0$ כך שלכל $k > k_0$ מתקיים $\sof{a_k}\frac{1}{k} < \frac{q}{\sof x}$. 
		
		ואז, לכל $k > k_0$ מתקיים $\sof{a_kx^{k}} < q^{k}$. כלומר המ''מ איברי הטור ''נשלטים`` ע''י האיברים של טור מתכנס, כלומר מבוחן $M$ של וויראשטראס הטור $\sum a_k x^{k}$ מתכנס בתחום המדובר. 
		
		נניח $\sof x > \bar R$ אז קיים $Q > 1$ כך ש־$\sof x > Q \bar R$. מהגדרת $\bar R$ נקבל:
		\[ \limsup _{k \to \infty}\sof{a_k}^{\frac{1}{k}} > \frac{Q}{\sof x} \]
		מההגדרה של $\limsup$ (תרגיל להראות זאת) יש ת''ס של $\sof{a_k}^{\frac{1}{k}}$ שהיא bounded away מ־$\frac{Q}{\sof x}$. ואז בת''ס $k_j$ הזו מתקיים: 
		\[ \sof{a_{k_j}x^{k_j}} > Q^{k_j} > 1 \]
		ומכאן שהטור $\sum a_k x^{k}$ אינו מתכנס בנקודה זו. 
	\end{proof}
	
	\textbf{דוגמאות. }
	\begin{itemize}
		\item במקרה של הטור $\sumkinf \frac{x^{k}}{k!}$ רדיוס ההתכנסות הוא $R = \infty$ והוא מתכנס ל־$e^{x}$. 
		\item במקרה של הטור $\sumkinf k!x^{k}$, קל לראות שהוא שואף לאינסוף בכל התום פרט ל־$x = 0$. ההוכחה המפורשת מקושי־הדמרד:
		\[ \cl{\frac{k}{2}}^{\frac{1}{2}} \le \cl{\cl{\ceil{\frac{k}{2}}}^{\ceil{\frac{k}{2}}}}^{\frac{1}{k}} \le \sqrt[k]{k!} \rrr{k \to \infty} \infty \]
		אז $R = 0$. 
		\item במקרה של הטור $\sumkinf \frac{x^{k}}{k}$ נקבל מקושי הדמרד רדיוס התכנסות $\limsup \sqrt[k]{\frac{1}{k}} = 1$. 
		\begin{itemize}
			\item ב־$x = 1$, זה טור הרמוני. הוא מתבדר. 
			\item ב־$x = -1$, ממשפט לייבניץ (חדו''א 1א) הוא מתכנס בתנאי. 
		\end{itemize}
		\item במקרה של הטור $\sumkinf (-1)^{k}\frac{x^{2k}}{k}$, רדיוס ההתכנסות $R = 1$ כי הגבול העליון: 
		\[ \limsup = \sqrt[m]{\sof{a_m}} = \lim\sqrt[2k]{\frac{1}{k}} = 1 \]
		ב־$x = \pm 1$ הטור מתכנס בתנאי. 
	\end{itemize}
	בשלב הזה יעל זורקת דוגמאות, הן לא ממש מעניינות אז לא אעתיק את כולן. בכללי כל החלק הזה בהרצאה לקוח מחדו''א 1א. 
	
	תרגיל: האם ייתכן התכנסות בהחלט בקצה אחד, והתבדרות בקצה השני? 
	
	\subsection*{גזירה איבר־איבר של טור חזקות}
	\theo{לטורים $\sumkoinf a_kx^{k}$ ו־$\sumkinf ka_kx^{k - 1}$ יש אותו רדיוס ההתכנסות. }\begin{proof}
		מחדו''א 1א ידוע ש־$\sqrt[k]{k} \to 1$. מכאן ש־: 
		\[ \limsup_{k \to \infty}\sqrt[k]{k \sof{a_k}} = \limsup_{k \to \infty} \sqrt[k]{\sof{a_k}} \]
		ומימין ניתן רדיוס ההתכנסות של הטור הראשון. 
	\end{proof}
	
	
	
	\ndoc
\end{document}